% !TEX root = ../main.tex
\chapter{Compilação}
\label{cap:compilacao}

Com o \code{media_movel.cmm} escrito, é hora de transformá-lo em hardware. Este capítulo explica o que acontece quando cada botão de compilação é acionado e quais arquivos aparecem no projeto.

\section{O pipeline do YANC em três estágios}
\label{sec:pipeline-yanc}

Todo clique em \botao{Compilar \cmm{}} executa, por processador, a cadeia de três compiladores do YANC. O \tool{cmmcomp} traduz o \file{<proc>.cmm} em \textit{assembly} simbólico, o \file{<proc>.asm}, junto com os registros e as tabelas que ligam cada instrução à linha do fonte. O \tool{appcomp} faz a primeira passada sobre o \textit{assembly}, coletando os parâmetros do processador e resolvendo os endereços de variáveis e rótulos. O \tool{asmcomp}, na segunda passada, gera o processador em Verilog (\file{Hardware/<proc>.v}), as duas imagens de memória, \file{<proc>_inst.mif} com o programa e \file{<proc>_data.mif} com os dados, e o \textit{testbench} (\file{Simulation/<proc>_tb.v}).

O terminal TCMM mostra a saída do primeiro estágio; o TASM, a dos outros dois, incluindo os avisos de recurso instanciado que revelam o custo de hardware do código (Seção~\ref{sec:custo-hardware}).

\begin{info}
O processador gerado instancia os moldes de HDL do SAPHO que acompanham a instalação: o núcleo, a ULA, as memórias. O programa não é um arquivo carregado depois; as imagens \file{.mif} nascem embutidas no projeto de hardware, e reprogramar significa recompilar e sintetizar de novo.
\end{info}

\section{Os botões de compilação}

O botão \botao{Compilar \cmm{}} roda o pipeline completo da seção anterior para os processadores do projeto, do fonte ao Verilog, memórias e \textit{testbench}. O \botao{Sintetizar Verilog} verifica os arquivos Verilog com o Icarus Verilog, apontando erros de sintaxe e elaboração no terminal TVERI, com as linhas clicáveis. O \botao{Analisar Verilog} percorre o fluxo completo até as formas de onda: compila o que for preciso, roda a simulação e abre o visualizador (Capítulo~\ref{cap:simulacao}). O \botao{Fast Sim} simula sem gravar ondas, disponível com o Verilator selecionado, para quando só interessam as saídas em arquivo. O \botao{PRISM} sintetiza com o Yosys e abre o diagrama de RTL (Capítulo~\ref{cap:prism}), exigindo um \textit{top-level} definido. O \botao{Teste do processador sintetizado} constrói um \textit{harness} C++ com o Verilator em torno do \file{<proc>.v} e exercita apenas as portas de entrada e saída, o teste de fumaça do processador como caixa-preta, com progresso no THTEST. E o \botao{Cancelar} interrompe todos os processos de compilação e simulação em andamento, sem fechar a IDE.

O botão \botao{Iniciar Compilação}, o raio da barra de status, é um atalho para o fluxo principal, e todos esses comandos também existem na paleta (\tecla{Ctrl+Shift+P}), no grupo de compilação.

Antes de compilar, a AURORA salva os arquivos automaticamente. O fluxo de ondas exige um \textit{testbench}, papel que o \file{<proc>_tb.v} gerado cumpre, e o PRISM exige um \textit{top-level}; a barra de status avisa o que falta, e o processador ativo acompanha o \file{.cmm} em foco no editor.

\section{Onde cada artefato aparece}

Depois de compilar o \code{media_movel}, o projeto ganha esta forma:

\begin{lstlisting}
media_movel/
|-- Software/
|   |-- media_movel.cmm        (seu fonte)
|   |-- media_movel.asm        (assembly gerado)
|   `-- pc_media_movel_mem.txt (tabela PC -> linha, p/ ondas)
|-- Hardware/
|   |-- media_movel.v          (o processador em Verilog)
|   |-- media_movel_inst.mif   (imagem da memoria de programa)
|   `-- media_movel_data.mif   (imagem da memoria de dados)
`-- Simulation/
    `-- media_movel_tb.v       (testbench gerado)
\end{lstlisting}

O \file{.v} e os \file{.mif} da pasta \file{Hardware} são exatamente o que se leva ao Quartus ou ao Vivado na hora de sintetizar para o FPGA real. O \textit{testbench} gerado produz \textit{clock} e \textit{reset}, alimenta o processador com os estímulos de \file{input_<i>.txt}, um valor por linha e um arquivo por porta de entrada, e grava as saídas em \file{output_<i>.txt}, além do despejo de ondas.

\begin{nota}
Segurança de execução: a AURORA só executa binários da própria cadeia de ferramentas, validados contra uma lista fechada, e os botões não passam por um \textit{shell}. É por isso que a IDE não oferece a execução de comandos arbitrários fora do terminal TCMD.
\end{nota}
